\chapter{风险、催化与监控}

本报告的核心风险不在于“预告会不会变差”，而在于市场将短期利润、周期价差或项目结算误读为可持续盈利。监控应围绕能改变分母、倍数和外部锚的变量，而不是围绕新闻标题。

\noindent
\begin{exhibitbox}[图表 10｜风险监控框架：从风险描述到可执行触发器]
\small
\begin{tabularx}{\textwidth}{L{2.2cm}L{4.0cm}L{4.4cm}X}
\toprule
\textbf{风险类别} & \textbf{领先指标} & \textbf{受影响模型} & \textbf{行动规则} \\
\midrule
盈利质量 & 扣非与归母差额扩大、非经常性收益上升、H2利润低于桥接。 & 常规EPS/PE、成长PE与恢复模型。 & 下调分母或切换至恢复模型；不以提高PE抵消盈利下修。 \\
现金流与营运资本 & 经营现金流持续低于利润、应收/库存占用上升。 & 周期、成长和项目型公司。 & 暂停升级；将牛市情景权重下调或取消。 \\
周期与价格 & 商品价格、炼化价差、运价、锂价、利用率反转。 & 周期调整PE与恢复模型。 & 先下调H2转化率，再下调PE区间；不直接沿用H1盈利。 \\
估值与拥挤度 & 当前价高于概率值、外部目标价陈旧、短期涨幅过大。 & 全部条件性模型。 & 转为观察；等待价格回撤或外部锚更新。 \\
证据质量 & 原始研报不可得、仅有摘要/转载、数据时点过旧。 & 外部锚和正式模型。 & 外部权重归零；仅保留本机构条件性区间。 \\
\bottomrule
\end{tabularx}
\normalsize
\end{exhibitbox}

\section{下一轮催化日历}

\noindent
\begin{tabularx}{\textwidth}{L{2.4cm}L{4.5cm}X}
\toprule
\textbf{时点} & \textbf{需要确认的变量} & \textbf{模型含义} \\
\midrule
下一次业绩披露 & H2扣非利润、毛利率、现金流和非经常性收益。 & 决定H2/H1转换率是否成立，影响全部EPS类模型。 \\
行业价格与订单更新 & 金属、能源、化工价差、运价、锂价、订单和交付。 & 决定周期模型的熊/基准/牛倍数与分母可持续性。 \\
当前研报或共识更新 & 2026E/2027E收入、EPS、目标价及评级变化。 & 决定外部锚是否可新增正权重，条件性模型是否可升级。 \\
资金与价格确认 & 行业日/周资金、价格位置、成交和回撤承接。 & 只用于确认市场是否接受基本面，不替代基本面证据。 \\
\bottomrule
\end{tabularx}

\noindent
\begin{sourcequalitybox}[数据边界]
候选Q1财务覆盖142/142，研报PDF覆盖137/142。行业日/周资金均为31/31完整表，但观察日分别截至7月14日和7月10日；连续流入公开表仅30/112，只作为辅助证据。
\end{sourcequalitybox}
